
\chapter{Asymptotic Analysis of Some Sequences}

In this section, we analyze the asymptotic behavior of some recurrences that arise in the analysis of regular expressions.

\begin{proposition}
  \label{prop:recurrence-sum}
  Consider the recurrence
  \begin{equation}
    x_{n+2} = \dfrac{a}{n}\sum_{i=1}^{n}x_i + b, \quad n \ge 1, \quad x_1=x_2=0.
  \end{equation}
  \begin{itemize}
    \item If $a<1$, then $x_n \sim \dfrac{b}{1-a}$.
    \item If $a=1$, then $x_n \sim b\ln n$.
    \item If $a>1$, then $x_n \sim \dfrac{b}{(a-1)\Gamma(a)} n^{a-1}$.
  \end{itemize}
\end{proposition}

\begin{proof}
  Consider the generating function $X(z) = \sum_{n=1}^{\infty} x_n z^n$ with the convention that $x_0=0$. We rewrite the recurrence as
  $$(n+2)x_{n+2} - 2x_{n+2} = a\sum_{i=1}^{n}x_i + bn.$$
  We multiply both sides by $z^{n+2}$ and sum over $n \ge 1$. The left-hand side becomes
  \begin{align*}
    \sum_{n=1}^{\infty}((n+2)x_{n+2} - 2x_{n+2})z^{n+2}
     & = \sum_{n=3}^{\infty}n x_{n}z^{n} - 2\sum_{n=3}^{\infty}x_{n}z^{n}     \\
     & = z \sum_{n=1}^{\infty}n x_{n}z^{n-1} - 2\sum_{n=1}^{\infty}x_{n}z^{n} \\
     & = z X'(z) - 2X(z).
  \end{align*}
  The right-hand side becomes
  \begin{align*}
    \sum_{n=1}^{\infty}(a\sum_{i=1}^{n}x_i + bn)z^{n+2}
     & = az^2\sum_{n=1}^{\infty}\sum_{i=1}^{n}x_i z^{n} + bz^2\sum_{n=1}^{\infty}nz^{n} \\
     & = az^2 \dfrac{X(z)}{1-z} + \dfrac{bz^3}{(1-z)^2}.
  \end{align*}

  Therefore, the generating function $X(z)$ satisfies the following ODE
  $$X'(z) - \left(\dfrac{2}{z} + \dfrac{az}{1-z}\right)X(z) = \dfrac{bz^3}{(1-z)^2}.$$

  We have $\mu(z) = \exp\left( \int -\left( \frac{2}{z} + \frac{a}{1-z} - a \right) \,\mathrm{d}z \right) = z^{-2} (1-z)^a$. which gives
  $$z^{-2} (1-z)^a X(z) = \int b (1-z)^{a-2} \,\mathrm{d}z.$$

  There are two cases of $a$ that determine the integral. If $a=1$, we have
  $$z^{-2} (1-z) X(z) = b\ln\left(\dfrac{1}{1-z}\right) + C.$$
  Since $x_1=x_2=0$, we have $z^{-2}X(z)\to 0$ as $z\to 0$, which gives $C=0$. Therefore, we have
  $$X(z) = b \dfrac{z^2}{1-z} \ln\left(\dfrac{1}{1-z}\right).$$
  Hence, $x_n = bH_{n-2} \sim b\ln n$.

  If $a \neq 1$, we have
  $$z^{-2} (1-z)^a X(z) = -\dfrac{b}{a-1} (1-z)^{a-1} + C.$$
  Using the same argument as above, we have $C= \dfrac{b}{a-1}$. Therefore, we have
  $$X(z) = \dfrac{b}{a-1}z^2 \left((1-z)^{-a}-(1-z)^{-1}\right).$$
  Using the generalized binomial theorem $$(1-z)^{-a} = \sum_{n=0}^{\infty} \binom{n+r-1}{n} z^n,$$ we have
  $$x_n = \dfrac{b}{a-1}\left(\binom{n-3+a}{n-2} - 1\right) \sim \dfrac{b}{a-1}\left(\dfrac{n^{a-1}}{\Gamma(a)}-1\right).$$

  If $a<1$ then $x_n \sim \dfrac{b}{1-a}$. If $a>1$, then $x_n  \sim \dfrac{b}{(a-1)\Gamma(a)} n^{a-1}$.
\end{proof}

Next, we want to see if adding a difference term to the recurrence changes the asymptotic behavior.
Consider the recurrence
\begin{equation}
  x_{n+2} = \dfrac{a}{n}\sum_{i=1}^{n}x_i + \dfrac{b}{n}\sum_{i=1}^{n}|x_i - x_{n+1-i}| + c, \quad n \ge 1, \quad x_1=x_2=0. \label{eq:recurrence-diff}
\end{equation}
where $a>1$, $0 < b < 1$ and $c>0$ are constants.


\begin{lemma}
  The sequence $(x_n)$ is increasing.
\end{lemma}

\begin{proof}
  One checks that $x_n\ge 0$ for all $n$. Let us prove that $(x_n)$ is increasing. We have $x_2\ge x_1$. Assume that $x_n\ge x_{n-1} \ge \ldots \ge x_0$ for some $n\ge 3$. Then,
  \begin{align*}
    n x_{n+2} & = a\sum_{i=1}^{n}x_i + b\sum_{i=1}^{n}|x_i - x_{n+1-i}| + nc                   \\
              & = a\sum_{i=1}^{n}x_i + b\sum_{i=1}^{\lfloor n/2\rfloor} (x_{n+1-i} - x_i) + nc \\
              & = a\sum_{i=1}^{n}x_i + b\sum_{i=1}^{\lfloor n/2\rfloor} (x_{n+1-i} - x_i) + nc \\
              & = a\sum_{i=1}^{n}x_i + b\sum_{i=1}^{\lfloor n/2\rfloor} (x_{n+1-i} - x_i) + nc \\
              & = a\sum_{i=1}^{n}x_i + b\sum_{i=1}^{\lfloor n/2\rfloor} (x_{n+1-i} - x_i) + nc \\
              & = (a+b)\sum_{i=1}^{n}x_i - 2b\sum_{i=1}^{\lfloor n/2\rfloor} x_i - bm_n + nc,
  \end{align*}

  where
  $$m_n =
    \begin{cases} x_{\lfloor n/2\rfloor+1} & \text{if } n \text{ is odd}, \\ 0 & \text{if } n \text{ is even}.
    \end{cases}$$

  Let $S_n = \sum_{i=1}^n x_i$. Then,
  \begin{align*}
    nx_{n+2} & = (a+b)S_n - 2b S_{\lfloor n/2\rfloor}  -bm_n + nc, \\
    S_{n+1}  & = S_n + x_{n+1}.
  \end{align*}

  Hence,
  $$nx_{n+2} - (n-1)x_{n+1} = (a+b)x_n - 2b(S_{\lfloor n/2\rfloor} - S_{\lfloor (n-1)/2\rfloor}) - b(m_n - m_{n-1}) + c.$$

  Let us check the right-hand side according to the parity of $n$.

  \begin{itemize}
    \item If $n=2k$, then $S_{\lfloor n/2\rfloor} - S_{\lfloor (n-1)/2\rfloor} = x_k$ and $m_n - m_{n-1} = -x_{k}$. The right-hand side is $(a+b)x_n - 2bx_k + c$.
    \item If $n=2k+1$, then $S_{\lfloor n/2\rfloor} - S_{\lfloor (n-1)/2\rfloor} = 0$ and $m_n - m_{n-1} = x_{k+1}$. The right-hand side is $(a+b)x_n - 2bx_{k+1} + c$.
  \end{itemize}

  Collecting the two cases, we have
  \begin{align*}
    nx_{n+2} - (n-1)x_{n+1} = (a+b)x_n - 2b x_{\lceil n/2\rceil} + c.
  \end{align*}
  The proof follows from Proposition \ref{prop:recurrence-ceil}.
\end{proof}

\begin{lemma}
  The sequence can be rewritten as
  \begin{equation}
    nx_{n+2} - (n-1)x_{n+1} = (a+b)x_n - 2b x_{\lceil n/2\rceil} + c.
  \end{equation}
\end{lemma}



